\subsection{FOL feasibility and ILP formulation reference}\label{lp:catalog}
Each entry fixes a problem variant, gives logical feasibility conditions,
and then a linear objective and integer constraints. Costs, profits, and
capacities are nonnegative unless stated otherwise. All graphs are finite;
undirected graphs are simple. A binary variable is the indicator of its
corresponding selection predicate. Quantifiers range over the indicated
input sets. The objective is separate from the FOL feasibility formula.

\begin{center}
\small\setlength{\tabcolsep}{3pt}
\begin{tabularx}{\linewidth}{@{}Y{1.8}Y{0.2}@{}}
\toprule
\textbf{Formulation} & \textbf{Page}\\\midrule
Metric TSP & \pageref{form:tsp}\\
Single-machine scheduling & \pageref{form:single}\\
Parallel-machine scheduling & \pageref{form:parallel}\\
Bin packing & \pageref{form:bin}\\
Set packing & \pageref{form:setpack}\\
Prize-collecting Steiner tree & \pageref{form:pcst}\\
Balanced cut & \pageref{form:cut}\\
$k$-minimum spanning tree & \pageref{form:kmst}\\
Capacitated minimum spanning tree & \pageref{form:cmst}\\
Feedback vertex set & \pageref{form:fvs}\\
Shortest superstring & \pageref{form:string}\\
3-dimensional matching & \pageref{form:3dm}\\
Fair matching with group quotas & \pageref{form:fair}\\
Capacitated facility location & \pageref{form:facility}\\
Maximum independent set & \pageref{form:independent}\\
Capacitated knapsack & \pageref{form:knapsack}\\
Higher-dimensional knapsack & \pageref{form:multiknapsack}\\
\bottomrule
\end{tabularx}
\end{center}

\subsubsection*{How to read the logical notation}
\textbf{Candidate interpretations.} Predicates such as $P(i)$ (select item)
and $A(i,j)$ (assign $i$ to $j$) describe a candidate solution. The displayed
FOL formula tests that interpretation. Searching over interpretations is
an additional optimization step; existentially quantifying a predicate
would be \emph{second-order}, not ordinary FOL.

\textbf{Counting.} $\exists^{=k}z\in D:P(z)$ abbreviates ``exactly $k$'',
using $k$ distinct witnesses satisfying $P$ and asserting that every
$P$-element equals one of them. For $k=0$, it means $\forall z\in D:\neg P(z)$.
Likewise, $\exists^{\le k}$ forbids $k+1$ distinct witnesses, and
$\exists^{\ge k}$ requires $k$ distinct witnesses. All these are finite
FOL expansions for a fixed input integer $k$.

\textbf{Capacity.} For nonnegative weights $w_i$ and bound $B$, write
\[
 \operatorname{Cap}_{D,w,B}(P):=
 \bigwedge_{\substack{A\subseteq D\\\sum_{i\in A}w_i>B}}
       \left(\bigvee_{i\in A}\neg P(i)\right).
\]
It forbids selecting every member of any overweight set, and is equivalent
to $\sum_{i\in D}w_i\mathbf1_{P(i)}\le B$. In $\operatorname{Cap}(A(\cdot,j))$,
the dot is the quantified object and $j$ is fixed.
The conjunction over subsets is a \emph{finite input-indexed schema},
not a first-order variable ranging over arbitrary sets. Expanding it can
be exponential even when the resulting ILP has just one capacity inequality.

\textbf{Connectivity and trees.} Let $Z(v)$ select vertices and $X(e)$
select edges. Write $E(S)$ for edges internal to $S$, $\delta(S)$ for edges
with exactly one endpoint in $S$, and $\delta(v)=\delta(\{v\})$.
Require selected edges to have selected endpoints:
\[
 \operatorname{Ends}(Z,X):=
 \bigwedge_{uv\in E}\bigl(X(uv)\Rightarrow Z(u)\land Z(v)\bigr).
\]
For each fixed $\varnothing\ne S\subsetneq V$, connectivity requires
\[
 \begin{aligned}
 &(\exists u\in S:Z(u))\land(\exists v\notin S:Z(v))\\
 &\hspace{12mm}\Rightarrow\bigvee_{e\in\delta(S)}X(e).
 \end{aligned}
\]
Call the conjunction of these formulas $\operatorname{Conn}(Z,X)$.
For the input graph's simple cycles $\mathcal C(G)$, define
\[
 \operatorname{Acyc}(X):=
 \bigwedge_{C\in\mathcal C(G)}\bigvee_{e\in E(C)}\neg X(e).
\]
Then $\operatorname{Tree}(Z,X)$ means Ends, Conn, Acyc, and
$\exists v:Z(v)$. Here too, the cut/cycle lists are fixed finite schemas;
these abbreviations do not claim that general reachability has a uniform
ordinary-FOL definition.

\paragraph{Reusable exact tree ILP.}\label{form:tree-template}
With $x_e,y_v\in\{0,1\}$, let $\mathcal T(x,y)$ denote
\[
\begin{aligned}
 x_{uv}&\le y_u,\quad x_{uv}\le y_v &&(uv\in E),\\
 \sum_{e\in E(S)}x_e&\le |S|-1 &&(S\subseteq V,\ |S|\ge2),\\
 \sum_{e\in E}x_e&=\sum_{v\in V}y_v-1,\qquad
 \sum_{v\in V}y_v\ge1.
\end{aligned}
\]
The subset inequalities make the selected graph a forest; a forest on
$q$ selected vertices has $q-c$ edges for $c$ components. The edge count
therefore forces $c=1$. This is exact, including singleton trees, but has
exponentially many subset constraints. Merely imposing $q-1$ edges is
insufficient: a cycle plus an isolated selected vertex can satisfy that count.

\subsubsection{Metric TSP}\label{form:tsp}
\textbf{Variant.} A complete undirected graph on $n\ge3$ cities, with
symmetric costs satisfying $c_{uw}\le c_{uv}+c_{vw}$. Select a minimum-cost
Hamiltonian cycle. Let $X(uv)$ mean that the tour uses edge $uv$.

\textbf{FOL.} With $Z(v)$ always true,
\[
 \bigl(\forall v:\exists^{=2}u\ne v:X(uv)\bigr)
 \land\operatorname{Conn}(Z,X).
\]
\textbf{ILP.} Binary $x_e$; minimize $\sum_e c_ex_e$, subject to
\[
\begin{aligned}
 \sum_{e\in\delta(v)}x_e&=2 &&(v\in V),\\
 \sum_{e\in\delta(S)}x_e&\ge2 &&(\varnothing\ne S\subsetneq V).
\end{aligned}
\]
Degree 2 gives a collection of cycles; the cut constraints force a single
tour. Triangle inequality supports the approximation algorithms, while
the same integer formulation also models nonmetric TSP.

\subsubsection{Single-machine scheduling}\label{form:single}
\textbf{Variant.} Jobs $J$, positive integer processing times $p_i$, weights
$w_i\ge0$, release time 0, no preemption or setup times. Minimize total
weighted completion time $\sum_iw_iC_i$. Put $H=\sum_i p_i$.

\textbf{FOL.} A candidate gives integer start times $S_i$ in the finite
time domain $\{0,\ldots,H\}$. Using the fixed arithmetic relations,
\[
\begin{aligned}
 &\forall i:\quad 0\le S_i,\quad S_i+p_i\le H,\\
 &\forall i\ne j:\quad
 (S_i+p_i\le S_j)\lor(S_j+p_j\le S_i).
\end{aligned}
\]
\textbf{ILP.} For each $i<j$, binary $z_{ij}=1$ means $i$ precedes $j$.
Minimize $\sum_iw_i(S_i+p_i)$, subject to
\[
\begin{aligned}
 S_i+p_i&\le S_j+H(1-z_{ij}) &&(i<j),\\
 S_j+p_j&\le S_i+Hz_{ij} &&(i<j),\\
 0\le S_i&\le H-p_i,\quad S_i\in\mathbb Z &&(i\in J).
\end{aligned}
\]
The bound $H$ makes the unused ordering inequality redundant. Positive
processing times exclude cyclic precedences. An optimal schedule can have
no idle time, so this horizon and integer start times lose no optimal
solution. A pure makespan objective on this one-machine, all-jobs-available
variant would always have optimum $H$.

\subsubsection{Parallel-machine scheduling}\label{form:parallel}
\textbf{Variant.} $m$ identical machines, all jobs available at time 0,
positive integer $p_i$, no preemption, precedences, or setup times.
Minimize makespan $T$. Let $A(i,h)$ assign job $i$ to machine $h$.

\textbf{FOL.} With integer $0\le S_i$ and $S_i+p_i\le T\le H$,
\[
\begin{aligned}
 &\forall i:\exists!h:A(i,h),\\
 &\forall i\ne j,\,h:\quad A(i,h)\land A(j,h)\\
 &\hspace{8mm}\Rightarrow
 (S_i+p_i\le S_j)\lor(S_j+p_j\le S_i).
\end{aligned}
\]
\textbf{ILP.} Binary $x_{ih}$ and integer $0\le T\le H$:
\[
\begin{aligned}
 \min\quad &T,\\
 \sum_{h=1}^m x_{ih}&=1 &&(i\in J),\\
 \sum_{i\in J}p_ix_{ih}&\le T &&(h=1,\ldots,m).
\end{aligned}
\]
Start-time variables are unnecessary in this ILP: run each machine's assigned
jobs consecutively in any order. Its completion time equals its load.
For unrelated machines replace $p_i$ by $p_{ih}$ and choose a valid horizon;
release times or precedences require additional scheduling constraints.

\subsubsection{Bin packing}\label{form:bin}
\textbf{Variant.} Pack $n$ items of sizes $0\le a_i\le B$ into as few
capacity-$B$ bins as possible; at most $n$ bins are needed. $A(i,b)$ means
assignment and $U(b)$ means bin $b$ is used.

\textbf{FOL.}
\[
\begin{aligned}
 &\forall i:\exists!b:A(i,b),\\
 &\forall i,b:\quad A(i,b)\Rightarrow U(b),\\
 &\forall b:\quad\operatorname{Cap}_{J,a,B}(A(\cdot,b)).
\end{aligned}
\]
\textbf{ILP.} Binary $x_{ib},y_b$; minimize $\sum_{b=1}^n y_b$:
\[
\begin{aligned}
 \sum_{b=1}^n x_{ib}&=1 &&(i\in J),\\
 x_{ib}&\le y_b &&(i\in J,\ b=1,\ldots,n),\\
 \sum_i a_ix_{ib}&\le By_b &&(b=1,\ldots,n).
\end{aligned}
\]
The linking constraint also handles zero-size items correctly. Optional
symmetry breaking $y_b\ge y_{b+1}$ orders the used bins without changing OPT.

\subsubsection{Set packing}\label{form:setpack}
\textbf{Variant.} Given subsets $S_j\subseteq U$ and profits $w_j$, select
a pairwise disjoint subfamily of maximum total profit. Use $P(j)$ for selection.

\textbf{FOL.}
\[
 \forall j\ne k:\quad
 (P(j)\land P(k))\Rightarrow S_j\cap S_k=\varnothing.
\]
Here disjointness is a fixed input relation, equivalently
$\neg\exists u:(u\in S_j\land u\in S_k)$.

\textbf{ILP.} Binary $x_j$:
\[
 \max\sum_j w_jx_j,\qquad
 \sum_{j:u\in S_j}x_j\le1\quad(u\in U).
\]
Every element can belong to at most one chosen set. Taking $w_j=1$ gives
maximum-cardinality set packing.

\subsubsection{Prize-collecting Steiner tree}\label{form:pcst}
\textbf{Variant.} Given an undirected graph, edge costs $c_e\ge0$, and
penalties $\pi_v\ge0$, select a nonempty tree; pay $\pi_v$ for each omitted
vertex. Nonterminals can have penalty 0. Use $Z(v)$ and $X(e)$ for selected
vertices and edges.

\textbf{FOL.} $\operatorname{Tree}(Z,X)$; for a rooted variant also require
$Z(r)$. Mandatory terminals satisfy $Z(v)$ regardless of their penalty.

\textbf{ILP.} Binary $x_e,y_v$, with the full tree system
$\mathcal T(x,y)$ from page~\pageref{form:tree-template}:
\[
 \min\quad\sum_{e\in E}c_ex_e+\sum_{v\in V}\pi_v(1-y_v).
\]
Add $y_r=1$ only when a root is prescribed. Internal Steiner vertices count
as selected vertices even when their penalty is zero. Connectivity is
essential: otherwise isolated selected vertices avoid penalties for free.
If the variant requires collected prize at least $Q$, additionally impose
$\sum_v\pi_vy_v\ge Q$; that is an extra quota assumption.

\subsubsection{Balanced cut}\label{form:cut}
\textbf{Variant.} Minimum-weight bisection of an undirected graph: choose
one side of size $q=\lfloor |V|/2\rfloor$. Let $P(v)$ select that side and
$D(uv)$ indicate a crossing edge.

\textbf{FOL.}
\[
\begin{aligned}
 &\exists^{=q}v:P(v),\\
 &\forall uv\in E:\quad D(uv)\Leftrightarrow(P(u)\mathbin\oplus P(v)).
\end{aligned}
\]
\textbf{ILP.} Binary $x_v,d_{uv}$; minimize $\sum_{uv\in E}c_{uv}d_{uv}$:
\[
\begin{aligned}
 \sum_v x_v&=q,\\
 d_{uv}&\ge x_u-x_v,\quad d_{uv}\ge x_v-x_u,\\
 d_{uv}&\le x_u+x_v,\quad d_{uv}\le2-x_u-x_v
                         &&(uv\in E).
\end{aligned}
\]
The four inequalities enforce XOR at every feasible binary point.
For a balance interval $L\le |S|\le U$, replace the equality by
$L\le\sum_vx_v\le U$ and use the corresponding counting bounds in FOL.

\subsubsection{\texorpdfstring{$k$}{k}-minimum spanning tree}\label{form:kmst}
\textbf{Variant.} Minimum-cost tree spanning exactly $k$ selected vertices
of $G$, where $1\le k\le |V|$. Vertices not in the tree are omitted.

\textbf{FOL.}
\[
 \operatorname{Tree}(Z,X)\land(\exists^{=k}v:Z(v)).
\]
\textbf{ILP.} Binary $x_e,y_v$:
\[
 \min\sum_e c_ex_e,\qquad
 \mathcal T(x,y),\qquad \sum_v y_v=k.
\]
In particular $\sum_ex_e=k-1$. A rooted variant adds $y_r=1$.
Check the convention: ``$k$ nonroot vertices'' means $k+1$ total vertices.
For nonnegative costs, an at-least-$k$ tree can be pruned to $k$ vertices
without increasing cost (preserving a required root), so the optimum values
agree under those assumptions.

\subsubsection{Capacitated minimum spanning tree}\label{form:cmst}
\textbf{Variant.} Root $r$, positive integer demands $q_v$ for $v\ne r$,
$q_r=0$, and branch capacity $Q$. Every component after deleting $r$ from
the spanning tree must have total demand at most $Q$.

\textbf{FOL.} Require $\operatorname{Tree}(Z,X)$ and $\forall v:Z(v)$.
For every chosen root edge $rb$, require
$\operatorname{Cap}_{V\setminus\{r\},q,Q}(\operatorname{Branch}_b)$.
Here $\operatorname{Branch}_b(v)$ means a selected-edge path from $b$ to $v$
avoiding $r$; expand it as the finite disjunction of path lengths
$0,\ldots,|V|-1$, with first-order quantified intermediate vertices.

\textbf{ILP.} Replace edges by both arcs but omit arcs entering $r$.
Binary $x_{uv}$ selects an arc directed away from $r$; integer $f_{uv}\ge0$
carries demand. Put $D=\sum_{v\ne r}q_v$. Minimize
$\sum_{uv}c_{uv}x_{uv}$, subject to
\[
\begin{aligned}
 \sum_{u:uv}x_{uv}&=1 &&(v\ne r),\\
 \sum_{u:uv}f_{uv}-\sum_{w:vw}f_{vw}&=q_v &&(v\ne r),\\
 \sum_{w:rw}f_{rw}&=D,\\
 0\le f_{uv}&\le Qx_{uv} &&(uv\text{ an arc}).
\end{aligned}
\]
Positive demands exclude disconnected directed cycles; one parent per
nonroot vertex then gives a rooted tree. Each root arc carries exactly its
branch demand. If zero demands are allowed, add connectivity/acyclicity
constraints explicitly. Unit demands give the branch-size version.

\subsubsection{Feedback vertex set}\label{form:fvs}
\textbf{Variant.} Delete a minimum-cost vertex set so the remaining
undirected graph is acyclic. Let $D(v)$ indicate deletion and $w_v$ its cost.

\textbf{FOL.}
\[
 \bigwedge_{C\in\mathcal C(G)}\bigvee_{v\in V(C)}D(v).
\]
Every input cycle must contain a deleted vertex.

\textbf{ILP.} Binary $d_v$:
\[
 \min\sum_vw_vd_v,\qquad
 \sum_{v\in V(C)}d_v\ge1\quad(C\in\mathcal C(G)).
\]
These cycle inequalities can be exponential in number. For the directed
variant use directed cycles. Feedback vertex set differs from vertex cover:
a tree has feedback vertex set size 0, even though its edges may need a
nonempty vertex cover.

\subsubsection{Shortest superstring}\label{form:string}
\textbf{Variant.} Given nonempty strings $s_i$ over alphabet $\Sigma$, find
a shortest word containing each as a contiguous substring. Let
$\ell_i=|s_i|$ and $H=\sum_i\ell_i$; concatenation proves this length bound.

\textbf{FOL.} A candidate has length $L\in\{0,\ldots,H\}$ and a character
function $a$ on its positions. Require
\[
\begin{aligned}
 \forall i\;\exists t:\quad
 &1\le t,\quad t+\ell_i-1\le L,\\
 &\forall k\in\{1,\ldots,\ell_i\}:\quad a(t+k-1)=s_i[k].
\end{aligned}
\]
Minimize $L$. The position and alphabet domains are finite; arithmetic and
input characters are fixed relations/functions.

\textbf{ILP (direct position model).} Binary $u_p$ activates position $p$;
$z_{p,a}$ puts character $a$ there; $o_{it}$ chooses an occurrence of $s_i$
starting at $t$. For $p=1,\ldots,H$, minimize $\sum_pu_p$ with
\[
\begin{aligned}
 \sum_{a\in\Sigma}z_{p,a}&=u_p &&(p=1,\ldots,H),\\
 u_p&\ge u_{p+1} &&(p<H),\\
 \sum_{t=1}^{H-\ell_i+1}o_{it}&\ge1 &&(\text{each }i),\\
 o_{it}&\le z_{t+k-1,s_i[k]} &&(\text{each valid }i,t,k).
\end{aligned}
\]
Here $1\le t\le H-\ell_i+1$ and $1\le k\le\ell_i$ in the last line.
The active positions form one prefix; exactly one character occupies each.
Every chosen occurrence must match all its characters. Conversely, every
superstring of length at most $H$ gives these variables. Contained strings
and duplicate strings are allowed by this model.

\subsubsection{3-dimensional matching}\label{form:3dm}
\textbf{Variant.} Disjoint coordinate sets $A,B,C$ and allowed triples
$\mathcal R\subseteq A\times B\times C$. Choose a maximum-weight collection
sharing no coordinate; $P(t)$ selects triple $t$.

\textbf{FOL.}
\[
 \forall u\in A\cup B\cup C:\quad
 \exists^{\le1}t\in\mathcal R:\bigl(P(t)\land u\in t\bigr).
\]
Membership $u\in t$ means that $u$ is one of the triple's coordinates.

\textbf{ILP.} Binary $x_t$:
\[
 \max\sum_{t\in\mathcal R}w_tx_t,\qquad
 \sum_{t\in\mathcal R:u\in t}x_t\le1\quad(u\in A\cup B\cup C).
\]
Set $w_t=1$ for maximum cardinality. In the \emph{perfect} version,
$|A|=|B|=|C|$ and every coordinate must be used: change $\le1$ to $=1$
in both the counting formula and ILP and solve feasibility.

\subsubsection{Fair matching with group quotas}\label{form:fair}
\textbf{Variant.} Bipartite $G=(L\cup R,E)$, matching profits $w_e$, and
fixed groups $E_g\subseteq E$ with integer quotas $\ell_g,u_g$.
For a group of left vertices $L_g$, use $E_g=\{uv\in E:u\in L_g\}$;
then the group count equals the number of its matched vertices.

\textbf{FOL.} With $M(e)$ denoting selection,
\[
\begin{aligned}
 &\forall v:\exists^{\le1}e\in\delta(v):M(e),\\
 &\forall g:\quad
 (\exists^{\ge\ell_g}e\in E_g:M(e))\land
 (\exists^{\le u_g}e\in E_g:M(e)).
\end{aligned}
\]
\textbf{ILP.} Binary $x_e$:
\[
\begin{aligned}
 \max\quad &\sum_{e\in E}w_ex_e,\\
 \sum_{e\in\delta(v)}x_e&\le1 &&(v\in L\cup R),\\
 \ell_g\le\sum_{e\in E_g}x_e&\le u_g &&(\text{each group }g).
\end{aligned}
\]
Groups may overlap; quotas may make the model infeasible. Set $w_e=1$ for
maximum cardinality. Proportional quotas can instead use
$\alpha_g\sum_ex_e\le\sum_{e\in E_g}x_e\le\beta_g\sum_ex_e$,
which is still linear for fixed rational $\alpha_g,\beta_g$.
Extra quotas do not automatically preserve the ordinary bipartite matching
LP's integrality.

\subsubsection{Capacitated facility location}\label{form:facility}
\textbf{Variant.} Clients $I$ have demands $d_i>0$; facilities $J$ have
capacities $Q_j$ and opening costs $f_j$. Each client uses exactly one open
facility. Let $c_{ij}$ be the \emph{total} cost of serving client $i$ at $j$.

\textbf{FOL.} Assignment $A(i,j)$ and opening $O(j)$ satisfy
\[
\begin{aligned}
 &\forall i:\exists!j:A(i,j),\\
 &\forall i,j:\quad A(i,j)\Rightarrow O(j),\\
 &\forall j:\quad\operatorname{Cap}_{I,d,Q_j}(A(\cdot,j)).
\end{aligned}
\]
\textbf{ILP (unsplittable demand).} Binary $x_{ij},y_j$:
\[
\begin{aligned}
 \min\quad &\sum_jf_jy_j+\sum_{i,j}c_{ij}x_{ij},\\
 \sum_jx_{ij}&=1 &&(i\in I),\\
 x_{ij}&\le y_j &&(i\in I,\ j\in J),\\
 \sum_i d_ix_{ij}&\le Q_jy_j &&(j\in J).
\end{aligned}
\]
If shipping costs are \emph{per unit}, use $d_ic_{ij}$ in the objective.
For splittable demand, use shipment variables $a_{ij}\ge0$ with
$\sum_ja_{ij}=d_i$, $\sum_ia_{ij}\le Q_jy_j$, and $a_{ij}\le d_iy_j$;
continuous shipments give a MILP rather than a pure ILP.

\subsubsection{Maximum independent set}\label{form:independent}
\textbf{Variant.} Select a maximum-weight set of pairwise nonadjacent
vertices; use weights 1 for maximum cardinality.

\textbf{FOL.}
\[
 \forall u,v:\quad uv\in E\Rightarrow\neg(P(u)\land P(v)).
\]
\textbf{ILP.} Binary $x_v$:
\[
 \max\sum_vw_vx_v,\qquad x_u+x_v\le1\quad(uv\in E).
\]
The complement is a vertex cover, but approximation factors do not
transfer simply by complementing an approximate solution.

\subsubsection{Capacitated knapsack}\label{form:knapsack}
\textbf{Variant.} One knapsack of capacity $B$, indivisible items with
sizes $a_i\ge0$ and profits $v_i\ge0$, each available once.

\textbf{FOL.} $\operatorname{Cap}_{J,a,B}(P)$ for the selected-item predicate
$P(i)$. Maximize the sum of selected profits.

\textbf{ILP.}
\[
 \max\sum_i v_ix_i,\qquad
 \sum_i a_ix_i\le B,\qquad x_i\in\{0,1\}.
\]
For bounded multiplicities replace the domain by
$x_i\in\{0,1,\ldots,b_i\}$; the FOL version can select individual copies.
With unlimited copies use $x_i\in\mathbb Z_{\ge0}$ and check for
zero-size, positive-profit items, which make the objective unbounded.

\subsubsection{Higher-dimensional knapsack}\label{form:multiknapsack}
\textbf{Variant.} Each selected item simultaneously consumes $a_{ri}\ge0$
units of each resource $r=1,\ldots,d$, with capacity $B_r$. Each item is
available once; its profit is $v_i$.

\textbf{FOL.}
\[
 \bigwedge_{r=1}^d\operatorname{Cap}_{J,a_r,B_r}(P).
\]
\textbf{ILP.}
\[
 \max\sum_i v_ix_i,\qquad
 \sum_i a_{ri}x_i\le B_r\quad(r=1,\ldots,d),
 \qquad x_i\in\{0,1\}.
\]
There is one capacity constraint per resource and one shared selection
variable per item. This models multiple resource budgets, not geometric
placement of rectangles or boxes. For several separate knapsacks with
assignment choices, introduce $x_{ib}$ and require $\sum_bx_{ib}\le1$.

\paragraph{Model-checking reminders.}
Keep domains, quantifier ranges, objective direction, and the precise
variant with each formulation. Cut/cycle schemas may be exponential.
Connectivity, acyclicity, and capacity are separate obligations; degree or
edge-count constraints alone rarely establish all three. Dropping all
integrality restrictions gives a relaxation, not automatically an exact
polynomial-time algorithm for the original problem.

\paragraph{Further formulation references.}
The formulations above are written out for the stated variants. Useful
primary references are the SCIP modelling book's
\href{https://scipbook.readthedocs.io/en/latest/routing.html}{TSP formulations},
\href{https://github.com/joaopedroso/scipbook/blob/master/docs/bpp.rst}{bin-packing model},
and \href{https://scipbook.readthedocs.io/en/latest/flp.html}{facility-location models};
the \href{https://dimacs11.zib.de/downloads.html}{DIMACS prize-collecting Steiner definitions};
the \href{https://optimization-online.org/wp-content/uploads/2006/05/1404.pdf}{CMST branch-capacity definition};
and \href{https://arxiv.org/abs/2301.03862}{proportionally fair matching}.

